\begin{figure}[H]
\centering
\begin{tikzpicture}
\begin{axis}[
    width=0.98\linewidth,
    height=0.58\linewidth,
    ymode=log,
    xmin=0,
    xmax=10000000,
    ymin=1,
    ymax=1000000,
    xlabel={$n$},
    ylabel={$d(n)$},
    scaled x ticks=false,
    xtick={0,2000000,4000000,6000000,8000000,10000000},
    xticklabels={$0$,$2\cdot10^6$,$4\cdot10^6$,$6\cdot10^6$,$8\cdot10^6$,$10^7$},
    grid=both,
    major grid style={black!12},
    minor grid style={black!6},
    tick label style={font=\small},
    label style={font=\small},
    legend columns=2,
    legend style={draw=none, fill=none, font=\small, at={(0.5,1.02)}, anchor=south},
    legend cell align=left,
]
\addplot[
    only marks,
    mark=*,
    mark size=0.35pt,
    draw=none,
    fill=black,
    fill opacity=0.16,
] table[x=n,y=d] {figures/data/d_sample.dat};
\addlegendentry{sample}

\addplot[
    only marks,
    mark=*,
    mark size=1.7pt,
    draw=red!70!black,
    fill=red!70!black,
] table[x=n,y=d] {figures/data/d_outliers.dat};
\addlegendentry{top ten}
\end{axis}
\end{tikzpicture}
\caption{Reference search parameters $d(n)$ for $1\leq n\leq 10^7$ on a logarithmic scale.  The gray points are the deterministic sample $n\equiv 0\pmod {5000}$, and the red points mark the ten largest observed values of $d(n)$ in this range.}
\label{fig:points}
\end{figure}
